\chapter{État de l'art}
\label{chap:etat_de_lart}

\section{Gestion des incidents de sécurité et SIEM}

La gestion des incidents de sécurité informatique s'appuie sur un cycle
largement standardisé dans l'industrie~: préparation, détection et analyse,
confinement/éradication/récupération, puis retour d'expérience. Ce cycle est
notamment formalisé par le \textit{National Institute of Standards and
Technology} (NIST) dans son guide \textit{Computer Security Incident
Handling Guide} \cite{nist80061}, qui reste une référence de facto pour la
structuration des opérations de sécurité.

Les plateformes \textbf{SIEM} automatisent la partie \og détection et
analyse\fg{} en centralisant les journaux d'événements de multiples sources
(systèmes d'exploitation, pare-feux, applications) et en appliquant des
règles de corrélation pour générer des alertes. Elles constituent aujourd'hui
un composant central de tout \textbf{SOC}, mais elles ne résolvent pas, par
elles-mêmes, le problème du \textbf{volume} d'alertes à traiter.

\section{Wazuh comme plateforme SIEM/XDR open source}

\textbf{Wazuh} \cite{wazuh_docs} est une plateforme de sécurité open source
issue d'un \textit{fork} d'OSSEC, qui combine les fonctions d'un
\textit{Host-based Intrusion Detection System} (HIDS), d'un collecteur de
journaux et d'un moteur de règles de corrélation. Son architecture repose
sur trois composants principaux~:

\begin{itemize}
  \item le \textbf{Wazuh Manager}, qui reçoit les événements des agents,
        applique les règles de détection et gère la configuration~;
  \item le \textbf{Wazuh Indexer}, un moteur de recherche et d'indexation
        basé sur \textbf{OpenSearch} \cite{opensearch_docs} (un projet dérivé
        d'Elasticsearch), qui stocke les alertes et permet leur recherche
        performante~;
  \item le \textbf{Wazuh Dashboard}, une interface web de visualisation.
\end{itemize}

Ce projet s'appuie exclusivement sur le \textbf{Wazuh Indexer} comme source
de données~: les alertes y sont interrogées via son API REST compatible
OpenSearch, ce qui découple totalement notre plateforme du Wazuh Manager et
du Dashboard.

\section{Fatigue d'alerte et automatisation (SOAR)}

La littérature professionnelle en sécurité opérationnelle documente de
longue date le phénomène de \textit{alert fatigue}~: les enquêtes annuelles
menées auprès des équipes SOC (par exemple celles du SANS Institute) montrent
de manière récurrente qu'une fraction importante des alertes générées ne
sont jamais examinées faute de temps, et qu'une proportion significative des
alertes traitées s'avèrent être des faux positifs. Ce constat a motivé
l'émergence de la catégorie des outils \textbf{SOAR}
(\textit{Security Orchestration, Automation and Response}), dont des
solutions commerciales comme Splunk SOAR, Microsoft Sentinel ou IBM QRadar
SOAR sont des représentants connus. Ces outils automatisent des
\textit{playbooks} — des séquences d'actions prédéfinies (enrichissement,
notification, blocage) — mais restent généralement limités à des règles
déterministes~: ils ne \textit{raisonnent} pas sur le contexte d'une alerte
au sens où le ferait un analyste humain.

\section{Agents fondés sur les grands modèles de langage}

L'émergence des \textbf{grands modèles de langage} (\textit{Large Language
Models}, LLM) a ouvert la voie à une nouvelle génération d'automatisation,
fondée non plus sur des règles fixes mais sur un raisonnement itératif en
langage naturel combiné à l'appel d'outils externes. Deux travaux
fondateurs éclairent la conception retenue dans ce projet~:

\begin{itemize}
  \item Wei \textit{et al.} \cite{cot2022} montrent que inciter un modèle à
        expliciter ses étapes de raisonnement (\textit{chain-of-thought})
        avant de produire une réponse améliore significativement la qualité
        de ses conclusions sur des tâches complexes~;
  \item Yao \textit{et al.} \cite{react2022} proposent le paradigme
        \textbf{ReAct} (\textit{Reasoning + Acting}), dans lequel un modèle
        alterne entre raisonnement explicite et appels d'outils externes
        (recherche d'information, exécution d'actions), le résultat de
        chaque action venant enrichir le contexte du raisonnement suivant.
        Ce paradigme correspond très directement à la \textbf{boucle
        agentique} implémentée par l'\textit{AI Investigator} de ce projet
        (\cref{chap:implementation})~: le modèle observe une alerte,
        décide d'appeler un outil (recherche d'alertes corrélées,
        réputation d'une IP, etc.), lit le résultat, et itère jusqu'à
        pouvoir formuler un verdict.
\end{itemize}

Le référentiel \textbf{MITRE ATT\&CK} \cite{mitre_attack}, qui catalogue les
tactiques et techniques employées par les attaquants, constitue par ailleurs
une base de connaissance couramment mobilisée — bien qu'elle ne soit pas
directement intégrée dans la version actuelle du projet — pour qualifier et
prioriser les verdicts produits par ce type d'agent.

\section{Comparatif des approches existantes}

Le \cref{tab:comparatif} positionne \projectname{} par rapport aux grandes
catégories d'outils existants.

\begin{table}[H]
\centering
\caption{Comparatif entre les approches existantes et \projectname}
\label{tab:comparatif}
\begin{tabular}{@{}p{3.4cm}p{3.6cm}p{3.6cm}p{3.6cm}@{}}
\toprule
\textbf{Critère} & \textbf{SIEM seul (ex. Wazuh natif)} & \textbf{SOAR à
playbooks} & \textbf{\projectname} \\
\midrule
Détection d'alertes & Oui (règles) & Non (en amont) & S'appuie sur Wazuh \\
Triage automatique & Non & Oui (déterministe) & Oui (raisonnement IA) \\
Justification du verdict & Non & Limitée (log d'actions) & Explicite,
tracée par action \\
Adaptation à un cas non prévu & Non & Non (règles figées) & Oui
(raisonnement contextuel) \\
Interrogation en langage naturel & Non & Non & Oui (assistant
conversationnel) \\
Coût / complexité de mise en œuvre & Faible & Élevé (licences) & Faible
(open source + API/LLM local) \\
\bottomrule
\end{tabular}
\end{table}

\section{Synthèse et positionnement du projet}

L'état de l'art fait apparaître un espace encore peu occupé par les
solutions \textbf{open source} entre, d'un côté, les SIEM qui détectent mais
ne concluent pas, et de l'autre, les SOAR commerciaux, coûteux, dont
l'automatisation reste déterministe et peu adaptable. \projectname{} se
positionne précisément dans cet espace~: il conserve Wazuh comme moteur de
détection éprouvé, et y ajoute une couche de raisonnement fondée sur un
agent LLM inspiré du paradigme ReAct, avec une exigence forte de
\textbf{traçabilité} (chaque action de l'agent est enregistrée) et de
\textbf{sobriété} des accès aux données (l'agent n'a, par construction,
aucun accès aux données d'identité des utilisateurs — voir
\cref{chap:conception}).
